\chapter*{Abstract}
\addcontentsline{toc}{chapter}{Abstract}
% 150--250 words. Update when the model comparison of Chapter 6 is measured.
Applicants for a Schengen short-stay visa in Türkiye prepare a file of about a dozen documents
from requirements that are spread over several official pages and differ between destination
states. A missing document, an expired passport or a name that differs between two documents can
cost the fee and the appointment. General-purpose AI assistants answer such questions fluently,
but they do not show which official sentence a requirement comes from, whether it is still in
force, or whether the applicant's own documents agree with each other.

\project{} turns the official requirements into versioned data in which every checklist item and
every rule cites its source passage. It builds a checklist for the applicant's destination and
profile, stores the uploaded documents encrypted, and checks their key fields with a deterministic
rule engine: missing documents, documents that are too old, and disagreements between documents
and with the trip. Every finding names the rule, the documents and fields involved, and the
official source.

The fall prototype covers the Netherlands and Norway for employed and student applicants. On two
synthetic application files with 17 labelled errors, the rule engine reports every error and
nothing else, and gives no finding on the error-free files. A language model that reads the fields
from the documents is being connected, and the same model will be compared with and without the
rule engine. The spring semester extends this to a labelled synthetic data set, a comparison with
human reviewers and the monitoring of official pages.

\noindent\textbf{Keywords:} visa application, document verification, rule engine,
vision-language model, source-grounded information
